\section{The Martyrs and the Saints at the Threshold}

Precious in the sight of the Lord is the death of his saints (Ps.\ 115:15).\endnote{Ps.\ 115:15 (Douay-Rheims); Ps.\ 115:6 in the Clementine Vulgate (116:15 Hebrew numbering).} When the Church inscribes a martyr or confessor in her calendar, she ordinarily keeps not his bodily birth into this fallen world, but his \emph{dies natalis}---his birthday into heaven, the day on which his mortal course closed and his eternal life with God began.

This transfiguration of death begins with St.\ Paul. Writing in chains to the Philippians, the Apostle sets life and death in a balance no pagan philosophy could construct:

\begin{quote}
For to me, to live is Christ: and to die is gain. And if to live in the flesh: this is to me the fruit of labour. And what I shall choose I know not. But I am straitened between two: having a desire to be dissolved and to be with Christ, a thing by far the better. But to abide still in the flesh is needful for you. (Phil.\ 1:21--24)
\end{quote}

Near the end of his course, writing to Timothy as the sword drew near, Paul turns to liturgical sacrifice: ``For I am even now ready to be sacrificed: and the time of my dissolution is at hand. I have fought a good fight: I have finished my course: I have kept the faith. As to the rest, there is laid up for me a crown of justice which the Lord the just judge will render to me in that day'' (2 Tim.\ 4:6--8).

Preaching on Philippians in Antioch, St.\ John Chrysostom draws out Paul's exacting balance. Calling death ``gain'' means ``I shall more clearly be present with Him; so that my death is rather a coming to life; they who kill me will work on me no dreadful thing, they will only send me onward to my proper life.'' Yet calling life in the flesh ``the fruit of his work'' refutes every heretic who despises created life or dares self-destruction:

\begin{quote}
How great was his philosophy! How hath he both cast out the desire of the present life, and yet thrown no reproach upon it! For in that he saith, ``to die is gain,'' by this he hath cast out the desire, but in that he saith, ``to live in the flesh is the fruit of my work,'' here he shows that the present life also is needful, if we use it as need is, if we bear fruit\ldots{} ``It is good for me to depart and be with Christ,'' for even death is a thing indifferent; since death itself is no ill, but to be punished after death is an ill. Nor is death a good, but it is good after our departure ``to be with Christ.'' What follows death is either good or ill.\endnote{St.\ John Chrysostom, \emph{Homilies on Philippians}, Hom.\ 3, on Phil.\ 1:21--24 (NPNF, ser.\ 1, vol.\ 13).}
\end{quote}

From that measure Chrysostom overturns worldly mourning: ``Let us grieve for sinners, not only when dying, but also while living. Let us joy for the just, not only while living, but also when dead.'' Sinners are pitiable even in health because they stand at enmity with God; the righteous are blessed even in death, because ``be they here, or be they there, [they] are with the King; and there, in a higher and nearer degree, not through an entrance, or by faith, but `face to face'\,'' (1 Cor.\ 13:12).

Nowhere in Christian antiquity does that Pauline longing burn with a fiercer Eucharistic flame than in St.\ Ignatius of Antioch's letter to the Romans as he traveled in chains under ten Roman soldiers toward the arena. Already sentenced by imperial power, the aged bishop feared only that Roman Christians might obtain his pardon and rob him of his offering:

\begin{quote}
I write to the Churches, and impress on them all, that I shall willingly die for God, unless ye hinder me. I beseech of you not to show an unseasonable good-will towards me. Suffer me to become food for the wild beasts, through whose instrumentality it will be granted me to attain to God. I am the wheat of God, and let me be ground by the teeth of the wild beasts, that I may be found the pure bread of Christ.\endnote{St.\ Ignatius of Antioch, \emph{Epistle to the Romans}, ch.\ 4, shorter recension (ANF, vol.\ 1); see also ch.\ 5 (``Let fire and the cross\ldots{} come upon me: only let me attain to Jesus Christ'').}
\end{quote}

Ignatius reverses the world's vocabulary of life and death. ``When I suffer,'' he writes, ``I shall be the freed-man of Jesus, and shall rise again emancipated in Him'' (ch.\ 4). ``Pardon me, brethren: do not hinder me from living, do not wish to keep me in a state of death\ldots{} Suffer me to obtain pure light: when I have gone thither, I shall indeed be a man of God. Permit me to be an imitator of the passion of my God'' (ch.\ 6). And he names the inward spring of that courage:

\begin{quote}
For though I am alive while I write to you, yet I am eager to die. My love has been crucified, and there is no fire in me desiring to be fed; but there is within me a water that liveth and speaketh, saying to me inwardly, Come to the Father. I have no delight in corruptible food, nor in the pleasures of this life. I desire the bread of God, the heavenly bread, the bread of life, which is the flesh of Jesus Christ, the Son of God\ldots{} and I desire the drink of God, namely His blood, which is incorruptible love and eternal life.\endnote{St.\ Ignatius of Antioch, \emph{Epistle to the Romans}, chs.\ 6--7, shorter recension (ANF, vol.\ 1).}
\end{quote}

The martyr does not hate the flesh; he hungers for the Flesh of Christ, hearing the living water of the Spirit murmur within him, \emph{Come to the Father}.

Most Christians, however, die not in arenas but in ordinary rooms after years of labor, habit, and domestic care. Two Doctors of the Church, St.\ Francis de Sales and St.\ Alphonsus Maria de Liguori, forged a daily schooling for that ordinary threshold.

In the \emph{Introduction to a Devout Life}, written for laypeople amid households and business, St.\ Francis de Sales bids the soul place itself in God's presence and imagine lying on a deathbed without hope of recovery:

\begin{quote}
O my soul! thou shalt one day depart out of this body! but when shall the time be? Shall it be in winter or in summer? In the city or in the country? By day or by night? Shall it be suddenly or after due preparation? By sickness or by accident? Shalt thou have leisure to make thy confession? Shalt thou be assisted by thy spiritual father? Alas! of all this we know nothing; one thing only is certain: we shall die, and sooner than we imagine.\endnote{St.\ Francis de Sales, \emph{Introduction to a Devout Life}, pt.\ 1, ch.\ 13 (Fifth Meditation: On Death), Catholic ed.\ (New York: Fr.\ Pustet), pp.\ 32--34.}
\end{quote}

In that hour ``the world shall end for you\ldots{} it shall be reversed before your eyes,'' so that vain pleasures shrink to shadows while ``the sins which before seemed very small will appear as large as mountains, and your devotion very small'' (pp.\ 32--33). Bidding farewell to wealth, company, kindred, and its own body, the departed soul takes flight to the right hand or to the left along ``no other [way] than that which it began here in this world'' (p.\ 33). Yet Francis's gentleness never chills human affection; he purifies it into charity that outlives the grave---``O my dear friends and relations! pardon me if I love you no more, but with a holy friendship, which may last eternally''---and closes with immediate resolutions of amendment, prayer to our Lady and the saints, and a single aromatic phrase: ``Make a nosegay of myrrh'' (p.\ 34).

A century and a half later, St.\ Alphonsus Liguori brought missionary directness to the same truths in \emph{Preparation for Death}. Taking Hebrews 9:27 as his text in the Fourth Consideration, he dismantles every evasion:

\begin{quote}
The sentence of death has been written against all men: you are man; you must die. ``Our other goods and evils,'' says St.\ Augustine, ``are uncertain: death alone is certain.'' It is uncertain whether the infant that is just born shall be poor or rich, whether he shall have good or bad health, whether he shall die in youth or in old age. But it is certain that he shall die\ldots{} For me, who am now writing, and for you, who read this little book, has been decreed the day and the moment when I shall no longer write, and you shall no longer read.\endnote{St.\ Alphonsus Maria de Liguori, \emph{Preparation for Death; or, Considerations on the Eternal Maxims}, trans.\ a Catholic clergyman (Dublin: James Duffy, 1843), Fourth Consideration, First--Third Points, pp.\ 60--68, and Fifth Consideration, First Point, p.\ 70.}
\end{quote}

``As your name has been one day inserted in the registry of baptisms; so it shall be one day written in the records of the dead,'' Alphonsus warns (p.\ 63), while every portrait and garment of the departed echoes Ecclesiasticus 38:23: ``Yesterday for me; to-day for thee'' (p.\ 64). Because ``nothing is more certain than death; but nothing is more uncertain than the hour of death'' (p.\ 70, citing Idiota), postponing conversion is like laughing on the road to the scaffold. Yet keeping \emph{memento mori} (remember death) before the eyes breeds liberty and calm, not gloom: St.\ Charles Borromeo kept a skull on his desk, Cardinal Baronius engraved the words upon his ring, Blessed Juvenal Ancina pondered ``What you are, I was; and what I am, you shall be,'' and a holy hermit asked on his deathbed why he remained so cheerful replied: ``I have always kept death before my eyes; and therefore, now that it has arrived, I see nothing new'' (pp.\ 67--68). Every consideration in Alphonsus opens straight into childlike trust before the Crucifix: ``My beloved Redeemer, I would not dare to appear before you, did I not see you hanging on the cross, lacerated, despised, and lifeless, for the love of me\ldots{} Your wounds, your blood, your death are my hope'' (p.\ 65).
